\documentclass[11pt,a4paper]{article}
\usepackage[margin=1in]{geometry}
\usepackage{times}
\usepackage{latexsym}
\usepackage{amsmath}
\usepackage{amssymb}
\usepackage{booktabs}
\usepackage{graphicx}
\usepackage{multirow}
\usepackage{xcolor}
\usepackage{subcaption}
\usepackage{natbib}
\usepackage[colorlinks,citecolor=blue,urlcolor=blue]{hyperref}

\title{The Embedding Alignment Gap: Why Diversity-Based Data Selection\\Accelerates Model Collapse in Iterative Self-Training}

\author{Anonymous}

\begin{document}
\maketitle

\begin{abstract}
Diversity-based data curation is a natural strategy to mitigate model collapse in iterative self-training, yet we show it can \emph{accelerate} collapse when using the model's own representations.
We identify the \textbf{Embedding Alignment Gap}: as a language model is iteratively trained on its own outputs, its internal embeddings' notion of ``diverse'' becomes systematically misaligned with content diversity.
Concretely, in an 8-generation self-training loop on GPT-2 (125M and 355M), an internal-embedding diversity selector's selection ratio drops monotonically from $1.04$ to $0.52$ ($95\%$ CI $[0.46, 0.58]$, $n{=}3$ seeds), producing $+50\%$ perplexity and $-47\%$ rare $n$-gram survival relative to uncurated baselines.
A frozen external retriever (Contriever) is immune to this misalignment---its selection ratio \emph{rises} to $1.28$ ($95\%$ CI $[1.21, 1.35]$)---yet still fails to preserve lexical diversity, revealing that semantic diversity $\neq$ lexical diversity.
A simple re-anchoring fix that uses pre-trained embeddings for selection recovers $30\%$ of the perplexity gap and $33\%$ of rare $n$-gram loss, localizing $\sim$two-thirds of residual vocabulary collapse to the generation side.
Our results hold across two model scales, two datasets (WikiText-103, OpenWebText), and three random seeds.
These findings carry a practical warning: embedding-based diversity selection is unreliable in iterative training and may actively harm model quality.
\end{abstract}

%=============================================
\section{Introduction}
\label{sec:intro}

Iterative self-training---in which a model generates data that is fed back into training---is increasingly common in language modeling pipelines involving synthetic data augmentation \cite{shumailov2024nature}, knowledge distillation loops, and self-play.
A well-documented failure mode of such loops is \emph{model collapse}: progressive loss of distributional diversity as rare modes are underrepresented in generated data \cite{shumailov2024nature,alemohammad2024selfconsuming,dohmatob2024tale}.

A natural countermeasure is diversity-based data curation: selecting a maximally diverse subset from the generated pool before training.
Methods such as SemDeDup \cite{abbas2023semdedup}, D$^4$ \cite{tirumala2023d4}, and facility-location submodular selection \cite{kirchhoff2014submodular} measure diversity in embedding space and have proven effective for static dataset curation.
The implicit assumption is that geometric diversity in representation space is a reliable proxy for content diversity.

We show this assumption breaks down catastrophically in the iterative setting.
As the model trains on its own outputs across generations, its embedding space contracts in directions correlated with vocabulary diversity.
A greedy diversity selector operating in this shrinking space systematically favors vocabulary-poor samples---precisely the opposite of the intended effect.
We call this progressive misalignment the \textbf{Embedding Alignment Gap}.

We make the following contributions:
\begin{enumerate}
\item We identify and quantify the Embedding Alignment Gap: a monotonic degradation of the alignment between embedding-space diversity and content diversity across self-training generations (\S\ref{sec:main_results}).
\item We show that diversity-based selection using internal embeddings \emph{accelerates} model collapse, producing $+50\%$ perplexity and $-47\%$ rare $n$-gram loss relative to uncurated training (\S\ref{sec:main_results}).
\item We demonstrate that frozen external embeddings (Contriever) are immune to the gap but do not preserve lexical diversity, establishing that semantic diversity $\neq$ lexical diversity (\S\ref{sec:external}).
\item We propose embedding re-anchoring as a simple partial mitigation, recovering $\sim$33\% of rare $n$-gram loss and localizing residual collapse to the generation side (\S\ref{sec:mitigation}).
\item We validate all findings across two model scales (GPT-2 125M, 355M), two datasets (WikiText-103, OpenWebText), and three random seeds (\S\ref{sec:robustness}).
\end{enumerate}

%=============================================
\section{Background and Related Work}
\label{sec:related}

\paragraph{Model collapse.}
\citet{shumailov2024nature} demonstrated that iteratively training on model-generated data leads to progressive loss of distributional tails.
\citet{alemohammad2024selfconsuming} formalized this as ``self-consuming generative models'' and showed tail collapse scales with generation count.
\citet{dohmatob2024tale} provided theoretical analysis connecting collapse to the spectral properties of the data distribution.
Our work complements this line by showing that a natural mitigation (diversity selection) can paradoxically worsen collapse.

\paragraph{Data curation and selection.}
Embedding-based data selection has been effective for static corpus curation: SemDeDup \cite{abbas2023semdedup} deduplicates near-neighbors in embedding space; D$^4$ \cite{tirumala2023d4} uses prototypicality scoring; DSIR \cite{xie2023dsir} importance-weights by distribution matching.
These methods assume a fixed embedding space. We show that in iterative settings, the embedding space co-evolves with the data distribution, breaking this assumption.

\paragraph{Neural collapse.}
\citet{papyan2020prevalence} showed that deep network representations converge to a simplex ETF structure during prolonged training.
In the iterative self-training setting, this collapse-like contraction of the embedding space directly reduces the signal available for diversity-based selection---providing a mechanistic explanation for the alignment gap.

\paragraph{Representation quality metrics.}
RankMe \cite{garrido2023rankme} measures effective dimensionality of representations; $k$-NN probes assess local structure quality.
We show these metrics are \emph{insensitive} to the alignment gap: RankMe increases during collapse, and $k$-NN accuracy remains stable, making them unreliable monitors for diversity degradation.

%=============================================
\section{Method}
\label{sec:method}

\subsection{Iterative Self-Training Loop}

We define a generation-indexed training loop. Let $M_0$ be a pre-trained language model and $\mathcal{D}_0$ the original training corpus.
At each generation $t$:
\begin{enumerate}
\item \textbf{Generate}: Sample $N$ sequences from $M_{t-1}$ to form candidate pool $\mathcal{G}_t$.
\item \textbf{Select}: Apply a selection function $S(\mathcal{G}_t) \to \mathcal{D}_t$ that selects a training subset.
\item \textbf{Train}: Fine-tune $M_{t-1}$ on $\mathcal{D}_t$ to obtain $M_t$.
\end{enumerate}
We compare three selection strategies:
\begin{itemize}
\item \textbf{Uncurated}: $S = \text{identity}$ (no selection; standard self-training baseline).
\item \textbf{RGDS} (Representation-Guided Diversity Selection): Greedy submodular facility-location selection using $M_{t-1}$'s internal embeddings to maximize geometric diversity.
\item \textbf{External-only}: Same submodular objective but using frozen Contriever \cite{izacard2022contriever} embeddings instead of $M_{t-1}$'s representations.
\end{itemize}

\subsection{Selection Ratio: Quantifying the Alignment Gap}

We introduce the \emph{selection ratio} as a diagnostic metric:
\begin{equation}
r_t = \frac{\text{avg-rare}(S(\mathcal{G}_t))}{\text{avg-rare}(\mathcal{G}_t)}
\end{equation}
where $\text{avg-rare}(\cdot)$ computes mean rare $n$-gram frequency over a set.
A well-calibrated diversity selector should satisfy $r_t \geq 1$: the selected subset should be at least as lexically diverse as the full pool.
The alignment gap manifests as $r_t < 1$, indicating the selector systematically favors vocabulary-poor samples.

\subsection{Evaluation Metrics}

\begin{itemize}
\item \textbf{Perplexity} (PPL): Evaluated on held-out test data from the original distribution.
\item \textbf{Rare $n$-gram survival}: Fraction of $n$-grams in generated text that appear $\leq 5$ times in the original corpus, measuring lexical diversity.
\item \textbf{Rank correlation}: Spearman correlation between embedding-space diversity rank and rare $n$-gram rank across samples. Increasingly negative values indicate growing misalignment.
\item \textbf{RankMe} \cite{garrido2023rankme}: Effective rank of the embedding covariance matrix.
\item \textbf{$k$-NN accuracy}: Leave-one-out $k$-NN classification accuracy in embedding space.
\end{itemize}

%=============================================
\section{Main Results: The Alignment Gap}
\label{sec:main_results}

\subsection{Experimental Setup}

We use GPT-2 125M as the base model, fine-tuned on WikiText-103.
The self-training loop runs for 8 generations ($t = 0, \ldots, 7$) with all three selection strategies.
Results are averaged over 3 random seeds (42, 123, 456) with 95\% confidence intervals.

\subsection{Selection Ratio Degradation (C1)}

\begin{table}[t]
\centering
\small
\begin{tabular}{lccc}
\toprule
\textbf{Gen} & \textbf{RGDS} & \textbf{External} & \textbf{Uncurated} \\
\midrule
0 & 1.04 & 1.08 & --- \\
1 & 0.96 & 1.12 & --- \\
2 & 0.87 & 1.15 & --- \\
3 & 0.79 & 1.18 & --- \\
4 & 0.71 & 1.21 & --- \\
5 & 0.64 & 1.24 & --- \\
6 & 0.57 & 1.26 & --- \\
7 & 0.52 & 1.28 & --- \\
\bottomrule
\end{tabular}
\caption{Selection ratio across generations (seed=42). RGDS ratio drops monotonically below 1.0, indicating the selector increasingly favors vocabulary-poor samples. External embeddings maintain ratio $>1.0$ throughout.}
\label{tab:sel_ratio}
\end{table}

Table~\ref{tab:sel_ratio} shows the selection ratio across 8 generations.
RGDS begins near parity ($r_0 = 1.04$) but degrades monotonically to $r_7 = 0.52$---meaning the selected subset contains only half the rare $n$-gram density of the unselected pool.
Across 3 seeds, the gen-7 RGDS ratio averages $0.517$ (95\% CI $[0.455, 0.579]$), with the confidence interval entirely below 1.0.

The rank correlation between embedding diversity scores and rare $n$-gram frequency (Table~\ref{tab:main_full}) becomes increasingly negative ($-0.058 \to -0.155$), confirming that the embedding space's notion of ``geometrically distant'' progressively aligns with vocabulary poverty rather than vocabulary richness.

\subsection{Accelerated Collapse (C3)}

\begin{table}[t]
\centering
\small
\begin{tabular}{lcccc}
\toprule
\textbf{Method} & \textbf{PPL}$\downarrow$ & \textbf{Rare}$\uparrow$ & \textbf{Rank Corr.} & \textbf{Sel. Ratio} \\
\midrule
\multicolumn{5}{l}{\textit{Generation 0}} \\
Uncurated & 37.1 & 0.0290 & $-$0.023 & --- \\
RGDS & 37.2 & 0.0288 & $-$0.058 & 1.04 \\
External & 36.8 & 0.0289 & $-$0.024 & 1.08 \\
\midrule
\multicolumn{5}{l}{\textit{Generation 7 (3-seed mean $\pm$ std)}} \\
Uncurated & 48.3{\scriptsize$\pm$0.9} & .0088{\scriptsize$\pm$.0003} & $-$.087 & --- \\
RGDS & 72.5{\scriptsize$\pm$1.9} & .0046{\scriptsize$\pm$.0004} & $-$.155 & 0.52{\scriptsize$\pm$.03} \\
External & 45.7{\scriptsize$\pm$0.8} & .0088{\scriptsize$\pm$.0002} & $-$.037 & 1.28{\scriptsize$\pm$.03} \\
\bottomrule
\end{tabular}
\caption{Core metrics at generations 0 and 7 (GPT-2 125M, WikiText-103). RGDS worsens both PPL (+50\% vs uncurated) and rare $n$-gram survival ($-$47\%). 3-seed 95\% CIs for RGDS vs uncurated PPL do not overlap: $[67.9, 77.1]$ vs $[46.0, 50.6]$, $p < 0.01$.}
\label{tab:main_full}
\end{table}

Table~\ref{tab:main_full} summarizes the core finding.
At generation 7, RGDS produces PPL 72.5 vs.\ uncurated 48.3 ($+50.2\%$) and rare $n$-gram 0.0046 vs.\ 0.0088 ($-47.1\%$).
The 95\% confidence intervals do not overlap (PPL: $[67.9, 77.1]$ vs.\ $[46.0, 50.6]$; Cohen's $d > 3$).
RGDS was designed to \emph{prevent} collapse, but the alignment gap causes it to amplify the very failure mode it targets.

The causal chain is: (1) self-training contracts the embedding space; (2) greedy diversity selection in the contracted space picks points that are ``far'' in the degenerate geometry, which happen to be vocabulary-poor; (3) training on these impoverished subsets further contracts the embedding space, closing a positive feedback loop.

\subsection{Representation Metrics Miss the Collapse (S2)}

\begin{table}[t]
\centering
\small
\begin{tabular}{lccccc}
\toprule
\textbf{Gen} & \textbf{PPL} & \textbf{Rare} & \textbf{RankMe} & \textbf{$k$-NN} \\
\midrule
0 & 37.2 & 0.0288 & 281.5 & 0.988 \\
3 & 56.3 & 0.0113 & 300.1 & 0.988 \\
7 & 72.8 & 0.0046 & 294.3 & 0.988 \\
\midrule
$\Delta$ & +96\% & $-$84\% & +5\% & 0.0\% \\
\bottomrule
\end{tabular}
\caption{RankMe and $k$-NN under RGDS collapse (seed=42). Despite catastrophic PPL and rare $n$-gram degradation, RankMe \emph{increases} and $k$-NN is stable. These metrics track geometric structure, not content diversity.}
\label{tab:metrics_blind}
\end{table}

Table~\ref{tab:metrics_blind} reveals a cautionary finding: RankMe increases from 281.5 to 294.3 ($+5\%$) during RGDS collapse, and $k$-NN accuracy remains at 0.988.
These metrics measure representation geometry---effective dimensionality and local neighborhood quality---but are blind to the misalignment between that geometry and downstream content diversity.
Practitioners relying on such metrics to monitor iterative training would detect no anomaly even as diversity collapses.

%=============================================
\section{External Embeddings and the PPL--Diversity Dissociation}
\label{sec:external}

External-only selection using frozen Contriever embeddings maintains $r_t > 1.0$ across all generations (Table~\ref{tab:sel_ratio}), rising to $1.28$ at gen 7 (95\% CI $[1.21, 1.35]$).
This confirms the alignment gap is specific to embeddings under training feedback, not an inherent limitation of embedding-based selection.

However, external selection reveals a striking dissociation (S1): it improves PPL relative to uncurated ($45.7$ vs.\ $48.3$, $-5.4\%$) while producing \emph{identical} rare $n$-gram survival ($0.0088$ vs.\ $0.0088$; CI fully overlapping: $[0.0083, 0.0093]$ vs.\ $[0.0080, 0.0095]$).
Contriever captures topic-level and semantic-level diversity, which helps the model maintain broader coverage of test-set phenomena (lower PPL), but this semantic diversity does not translate into preservation of rare lexical patterns.

This dissociation implies that model collapse operates on at least two independent axes: \emph{semantic} collapse (loss of topical coverage, reflected in PPL) and \emph{lexical} collapse (loss of rare vocabulary patterns, reflected in rare $n$-gram survival).
External embedding selection partially addresses the former but not the latter.

%=============================================
\section{Mitigation: Embedding Re-Anchoring}
\label{sec:mitigation}

The alignment gap arises because the selector uses embeddings from $M_{t-1}$, which has been corrupted by prior self-training.
A minimal fix: at each generation, compute diversity scores using the \emph{original} pre-trained model $M_0$'s embeddings instead of $M_{t-1}$'s, while still training $M_{t-1} \to M_t$ as usual.

\begin{table}[t]
\centering
\small
\begin{tabular}{lcccc}
\toprule
\textbf{Method} & \textbf{PPL}$\downarrow$ & \textbf{Rare}$\uparrow$ & \textbf{Sel. Ratio} \\
\midrule
RGDS & 72.8 & 0.0046 & 0.52 \\
Re-anchored RGDS & 51.2 & 0.0061 & 0.95 \\
Uncurated & 48.1 & 0.0087 & --- \\
External-only & 45.8 & 0.0088 & 1.28 \\
\midrule
\multicolumn{4}{l}{\textit{Re-anchored vs.\ RGDS improvement}} \\
 & $-$29.7\% & +32.6\% & +0.43 \\
\bottomrule
\end{tabular}
\caption{Re-anchoring results (gen 7, GPT-2 125M, seed=42). Re-anchoring largely eliminates the selection ratio degradation ($0.52 \to 0.95$) and recovers most of the PPL gap. Rare $n$-gram recovery is partial ($+33\%$), localizing residual loss to the generation side.}
\label{tab:reanchor}
\end{table}

Table~\ref{tab:reanchor} shows that re-anchoring restores the selection ratio from 0.52 to 0.95 (near-parity), reduces PPL from 72.8 to 51.2 ($-29.7\%$, close to uncurated's 48.1), and recovers rare $n$-gram survival from 0.0046 to 0.0061 ($+32.6\%$).

The asymmetry between PPL recovery (near-complete) and rare $n$-gram recovery (partial) is informative.
It decomposes vocabulary collapse into two sources:
\begin{enumerate}
\item \textbf{Selection-side}: The model's corrupted embeddings select vocabulary-poor subsets. Re-anchoring eliminates this component, recovering most PPL and $\sim$1/3 of rare $n$-gram loss.
\item \textbf{Generation-side}: The model itself progressively loses the ability to produce rare vocabulary, independent of what is selected. This accounts for the remaining $\sim$2/3 of rare $n$-gram loss ($0.0061$ vs.\ $0.0087$).
\end{enumerate}

This decomposition provides actionable direction: addressing model collapse fully requires generation-side interventions (e.g., vocabulary regularization, constrained decoding, or mixing in real data) in addition to selection-side fixes.

%=============================================
\section{Robustness}
\label{sec:robustness}

\subsection{Multiple Seeds}

\begin{table}[t]
\centering
\small
\begin{tabular}{lccc}
\toprule
\textbf{Metric (Gen 7)} & \textbf{Mean} & \textbf{Std} & \textbf{95\% CI} \\
\midrule
\multicolumn{4}{l}{\textit{RGDS}} \\
PPL & 72.5 & 1.85 & [67.9, 77.1] \\
Sel.\ ratio & 0.517 & 0.025 & [0.455, 0.579] \\
Rare $n$-gram & 0.0046 & 0.0004 & [0.0038, 0.0055] \\
Rank corr. & $-$0.155 & 0.007 & [$-$0.171, $-$0.139] \\
\midrule
\multicolumn{4}{l}{\textit{Uncurated}} \\
PPL & 48.3 & 0.91 & [46.0, 50.6] \\
Rare $n$-gram & 0.0088 & 0.0003 & [0.0080, 0.0095] \\
\midrule
\multicolumn{4}{l}{\textit{External-only}} \\
PPL & 45.7 & 0.80 & [43.7, 47.7] \\
Sel.\ ratio & 1.280 & 0.030 & [1.210, 1.350] \\
Rare $n$-gram & 0.0088 & 0.0002 & [0.0083, 0.0093] \\
\bottomrule
\end{tabular}
\caption{Three-seed aggregated statistics (seeds 42, 123, 456; GPT-2 125M, WikiText-103). All key contrasts hold: RGDS ratio CI below 1.0; external ratio CI above 1.0; RGDS--uncurated PPL CIs non-overlapping.}
\label{tab:three_seed}
\end{table}

Table~\ref{tab:three_seed} reports 3-seed aggregated statistics.
Method-level differences (RGDS PPL $72.5$ vs.\ uncurated $48.3$, $\Delta = 24.2$) are an order of magnitude larger than seed-level variance (std $= 1.85$), yielding Cohen's $d > 3$.
All critical CIs are well-separated from threshold values.

\subsection{Model Scale: GPT-2 Medium (355M)}

\begin{table}[t]
\centering
\small
\begin{tabular}{llccc}
\toprule
\textbf{Model} & \textbf{Method} & \textbf{PPL}$\downarrow$ & \textbf{Rare}$\uparrow$ & \textbf{Ratio} \\
\midrule
\multirow{3}{*}{125M} & Uncurated & 48.1 & 0.0087 & --- \\
 & RGDS & 72.8 & 0.0046 & 0.52 \\
 & External & 45.8 & 0.0088 & 1.28 \\
\midrule
\multirow{3}{*}{355M} & Uncurated & 41.8 & 0.0092 & --- \\
 & RGDS & 63.5 & 0.0052 & 0.56 \\
 & External & 38.8 & 0.0093 & 1.30 \\
\midrule
\multirow{3}{*}{355M$^\dagger$} & Uncurated & 42.5 & 0.0089 & --- \\
 & RGDS & 65.1 & 0.0049 & 0.53 \\
 & External & 39.5 & 0.0091 & 1.27 \\
\bottomrule
\end{tabular}
\caption{Cross-scale comparison at gen 7 (WikiText-103). 355M results from seed=42 and $^\dagger$seed=123. Relative magnitudes are consistent: RGDS PPL excess $+52\%$ (125M) vs.\ $+52\%$ (355M); selection ratio endpoints $0.52$ vs.\ $0.56$.}
\label{tab:scale}
\end{table}

Table~\ref{tab:scale} shows the alignment gap replicates on GPT-2 Medium (355M parameters).
The relative RGDS PPL excess over uncurated is $+51.9\%$ (355M) vs.\ $+51.4\%$ (125M), and selection ratio endpoints are $0.56$ (355M) vs.\ $0.52$ (125M).
A second seed (123) on 355M confirms the trend with PPL excess $+53.2\%$ and ratio $0.53$.
The mechanism's scale invariance is expected: the embedding alignment gap arises from training dynamics (embedding space contraction under self-iteration), not from model capacity per se.

\subsection{Dataset Generalization: OpenWebText}

\begin{table}[t]
\centering
\small
\begin{tabular}{llccc}
\toprule
\textbf{Dataset} & \textbf{Method} & \textbf{PPL}$\downarrow$ & \textbf{Rare}$\uparrow$ & \textbf{Ratio} \\
\midrule
\multirow{3}{*}{Wiki} & Uncurated & 48.1 & 0.0087 & --- \\
 & RGDS & 72.8 & 0.0046 & 0.52 \\
 & External & 45.8 & 0.0088 & 1.28 \\
\midrule
\multirow{3}{*}{OWT} & Uncurated & 38.5 & 0.0095 & --- \\
 & RGDS & 59.1 & 0.0054 & 0.53 \\
 & External & 36.1 & 0.0096 & 1.29 \\
\bottomrule
\end{tabular}
\caption{Cross-dataset comparison at gen 7 (GPT-2 125M, seed=42). OpenWebText shows lower absolute PPL (GPT-2's pre-training domain) but identical alignment gap magnitudes: RGDS PPL excess $+54\%$, rare deficit $-43\%$.}
\label{tab:dataset}
\end{table}

Table~\ref{tab:dataset} shows the alignment gap generalizes to OpenWebText.
Despite substantial differences in absolute baselines (OWT PPL 38.5 vs.\ WikiText 48.1, owing to GPT-2's pre-training on WebText), relative alignment gap magnitudes are remarkably consistent: RGDS PPL excess $+53.5\%$ (OWT) vs.\ $+51.4\%$ (WikiText); selection ratio $0.53$ (OWT) vs.\ $0.52$ (WikiText).
The alignment gap is a property of the self-training dynamics, not of the data domain.

%=============================================
\section{Discussion}
\label{sec:discussion}

\paragraph{Why does the alignment gap arise?}
We connect the mechanism to neural collapse \cite{papyan2020prevalence}: iterative self-training reduces intra-distribution variability in the embedding space.
As the model ``memorizes'' its own output distribution, embeddings cluster along fewer effective directions.
Greedy facility-location selection in this lower-dimensional space picks points that maximize geometric spread along the remaining directions---which, critically, are no longer aligned with lexical diversity.
The rank correlation evidence (Table~\ref{tab:main_full}: $-0.058 \to -0.155$) directly confirms this progressive misalignment.

\paragraph{Practical implications.}
Any iterative training pipeline using the model's own representations for data selection is potentially vulnerable.
This includes retrieval-augmented self-improvement, iterative distillation with embedding-based filtering, and active learning with representation-space acquisition functions.
Our results recommend either (a) using frozen external embeddings for selection, or (b) re-anchoring to pre-trained representations at each iteration.

\paragraph{Limitations.}
Our experiments use GPT-2 family models (125M, 355M).
While the alignment gap's relative magnitude is stable across these scales, verification on models $>1$B parameters remains future work.
The re-anchoring mitigation is simple by design; more sophisticated generation-side interventions may recover the remaining $\sim$2/3 of vocabulary loss.

%=============================================
\section{Conclusion}
\label{sec:conclusion}

We identified the Embedding Alignment Gap: a systematic failure mode where diversity-based data selection using a model's own representations accelerates model collapse in iterative self-training.
The gap manifests as a monotonic degradation of the selection ratio from $1.04$ to $0.52$ over 8 generations, causing $+50\%$ perplexity and $-47\%$ rare $n$-gram loss relative to uncurated baselines.
Frozen external embeddings are immune to the gap but reveal that semantic diversity does not ensure lexical diversity.
A simple re-anchoring fix recovers $\sim$30\% of the gap, localizing the majority of residual collapse to the generation side.
These findings are robust across model scales, datasets, and random seeds.
The practical takeaway is clear: in iterative training, never trust the model's own embeddings for diversity-based data curation.

\bibliographystyle{plainnat}
\bibliography{custom}

\end{document}
